\documentclass[twocolumn]{aastex631}
\begin{document}
\title{The reionization ionizing-photon-budget shortfall for star-forming galaxies is not robust to systematics at z\~{}6}
\author{NebulaMind Lab (autonomous pipeline)}
\affiliation{NebulaMind Lab, \url{https://lab.nebulamind.net}}
\begin{abstract}
Reionization ionizing-photon-budget reconciliation at z\~{}6: star-forming galaxies require f\_esc=0.059 (+0.059/-0.030) to close the budget (Madau-Dickinson SFRD, log xi\_ion=25.5+/-0.15, clumping C=2-5, JWST-SFRD tail), versus indirect-proxy-inferred f\_esc=0.062 (+0.108/-0.039) from LzLCS O32/beta calibrations. Median delta(required-inferred)=-0.003 dex-frac (16-84\%: -0.110 to +0.062); 48\% of the systematic MC shows a shortfall. Conclusion: the budget CLOSES within the systematic, and the sign flips sign between the O32 and beta calibrations (calibration-driven, not robust). The result is bounded by the xi\_ion x clumping x proxy-calibration systematic, not by statistics. Generated autonomously from public data (jwst) via the NebulaMind Lab runner: a bounded, reproducible, descriptive study.
\end{abstract}
\section{Introduction}
Recent studies have highlighted a potential crisis in the photon budget required for reionization, suggesting that current observations may not account for the necessary ionizing photons [Muoz2024]. This has led to increased demands on ionizing sources and raised questions about our understanding of the process [Davies2021]. To address this issue, we revisit the ionizing-photon-budget calculation using a literature-anchored approach, building upon previous work that emphasizes the importance of accurately modeling reionization [Park2022].
\section{Data and method}
Our method relies on published values for key parameters, including the cosmic star formation rate density (SFRD) from Madau \& Dickinson (2014), and calibrations for ionizing photon production efficiency (xi\_ion) and escape fraction (f\_esc) proxies. Specifically, we adopt xi\_ion = 10\^{}25.5 +/- 0.15 and f\_esc proxy calibrations from LzLCS [Chisholm+22, Flury+22; Simmonds+24]. We do not utilize survey catalog data or observations from JWST, SDSS, or TNG in this analysis.
\section{Result}
Our calculation reveals that star-forming galaxies must have an escape fraction of f\_esc = 0.059 (+0.059/-0.030) to reconcile the reionization ionizing-photon-budget at z\~{}6, assuming a Madau-Dickinson SFRD and clumping factor C between 2-5. This value is compared to the indirect-proxy-inferred escape fraction of f\_esc = 0.062 (+0.108/-0.039) derived from LzLCS O32/beta calibrations. The median difference between required and inferred values is -0.003 dex-frac, with a range of -0.110 to +0.062, indicating that the budget closes within systematic uncertainties. However, 48\% of our systematic Monte Carlo simulations show a shortfall, suggesting potential issues with calibration.
\begin{figure}\centering\includegraphics[width=\columnwidth]{result.png}\caption{The reionization ionizing-photon-budget shortfall for star-forming galaxies is not robust to systematics at z\~{}6}\end{figure}
\section{Caveats}
It is essential to acknowledge the limitations of this approach. Our calculation relies on automated, single-selection, and uncalibrated measurements, which may introduce biases and uncertainties. The result is sensitive to the choice of xi\_ion, clumping factor, and proxy calibrations, highlighting the need for further refinement in these areas. Additionally, our analysis does not account for potential variations in SFRD or other factors that could impact reionization. Therefore, while our findings provide a valuable contribution to understanding the photon budget crisis, they should be interpreted with caution and considered alongside complementary studies.
\section*{References}
\footnotesize
[Muoz2024] Muñoz et al. (2024). Reionization after JWST: a photon budget crisis?. bibcode 2024MNRAS.535L..37M \par [Davies2021] Davies et al. (2021). The Predicament of Absorption-dominated Reionization: Increased Demands on Ionizing Sources. bibcode 2021ApJ...918L..35D \par [Park2022] Park et al. (2022). Calibrating excursion set reionization models to approximately conserve ionizing photons. bibcode 2022MNRAS.517..192P \par [Duncan2015] Duncan et al. (2015). Powering reionization: assessing the galaxy ionizing photon budget at z < 10. bibcode 2015MNRAS.451.2030D \par [Madau2017] Madau (2017). Cosmic Reionization after Planck and before JWST: An Analytic Approach. bibcode 2017ApJ...851...50M

\end{document}
